\section{A Gaussian Core Characterization}
\label{sec:gaussian}

Suppose a paired-oracle call to $v\in V(H)$ returns
\[
Z_v\sim\mathcal{N}(\theta_v,\nu_v^2),
\]
independently across paired-oracle calls, with known $\nu_v^2>0$.  We test
\[
\mathsf H_0:\theta_u=\theta_v\ \forall (u,v)\in E(H)
\]
against
\[
\mathsf H_1:\exists (u,v)\in E(H)
\text{ such that }|\theta_u-\theta_v|\ge\gamma.
\]

\begin{definition}[Uncertainty-weighted graph complexity]
\label{def:psi}
\[
\Psi(H,\nu)=
\min_{q_v>0}\sum_{v\in V(H)}q_v
\quad\text{s.t.}\quad
\frac{\nu_u^2}{q_u}+\frac{\nu_v^2}{q_v}\le1
\quad\forall (u,v)\in E(H).
\]
\end{definition}

The program is convex.  One paired-oracle call at a vertex supports every incident edge,
which is why the optimum depends on graph structure rather than simply $|E(H)|$.

\begin{theorem}[Static upper bound]
\label{thm:upper}
Let $m=|E(H)|\ge1$.  For $\delta\in(0,1)$, a non-adaptive tester exists using at
most
\[
\frac{8\log(2m/\delta)}{\gamma^2}\Psi(H,\nu)+|V(H)|
\]
paired-oracle calls and having error probability at most $\delta$.
\end{theorem}

\begin{proof}
Take an optimal solution $q^\star$ and set
\[
n_v=\left\lceil
\frac{8\log(2m/\delta)}{\gamma^2}q_v^\star
\right\rceil.
\]
For each edge $(u,v)$, the empirical difference of the two vertex means has
variance at most $\gamma^2/(8\log(2m/\delta))$.  Reject if an empirical edge
difference has magnitude at least $\gamma/2$.  Under the null, a Gaussian tail
bound followed by a union bound controls all edges.  Under the alternative, an
edge with true difference at least $\gamma$ can be missed only if its estimation
error is at least $\gamma/2$.  The ceilings contribute at most $|V(H)|$ paired-oracle calls.
\end{proof}

Let $\kl(a,b)$ denote binary relative entropy.

\begin{theorem}[Adaptive lower bound]
\label{thm:lower}
For $\delta\in(0,1/2)$, every adaptive $\delta$-correct tester with finite
expected stopping time (measured in paired-oracle calls) has, on the null instance $\theta\equiv0$,
\[
\E_0[T]
\ge
\frac{2\kl(1-\delta,\delta)}{\gamma^2}\Psi(H,\nu).
\]
\end{theorem}

\begin{proof}
Let $n_v=\E_0[N_v]$.  Fix an edge $(u,v)$ and consider an alternative changing
only its endpoints to $\theta_u=a$ and $\theta_v=-b$, where $a,b\ge0$ and
$a+b=\gamma$.  The adaptive transcript KL from the null to this alternative is
\[
\frac{n_u a^2}{2\nu_u^2}+\frac{n_v b^2}{2\nu_v^2}.
\]
The standard adaptive change-of-measure argument for sequential sampling
\cite{KaufmannCappeGarivier2016}, followed by data processing through the final
accept/reject decision, lower-bounds this by $\kl(1-\delta,\delta)$.  Minimizing over $a+b=\gamma$ yields
\[
\frac{\nu_u^2}{n_u}+\frac{\nu_v^2}{n_v}
\le
\frac{\gamma^2}{2\kl(1-\delta,\delta)}.
\]
This holds for every edge.  Scaling $n_v$ by
$\gamma^2/(2\kl(1-\delta,\delta))$ therefore gives a feasible solution to
\Cref{def:psi}, proving the claim.
\end{proof}

\subsection{Dual characterization}

\begin{theorem}[Dual representation]
\label{thm:dual}
\[
\Psi(H,\nu)
=
\max_{\lambda_e\ge0}
\left[
2\sum_{v\in V(H)}\nu_v
\sqrt{\sum_{e\ni v}\lambda_e}
-
\sum_{e\in E(H)}\lambda_e
\right].
\]
\end{theorem}

\begin{proof}
Introduce one multiplier $\lambda_e\ge0$ for each edge constraint.  The Lagrangian
separates across vertices as
\[
\sum_v\left(q_v+
\frac{\nu_v^2\sum_{e\ni v}\lambda_e}{q_v}\right)-\sum_e\lambda_e.
\]
For fixed multipliers, minimizing over $q_v>0$ gives
$2\nu_v\sqrt{\sum_{e\ni v}\lambda_e}$.  Strong duality follows from Slater's
condition because choosing all $q_v$ sufficiently large makes every inequality
strict.
\end{proof}

\subsection{Exact unit-variance families}
Assume $\nu_v^2=1$.

\begin{proposition}[Closed forms]
\label{prop:families}
\begin{enumerate}[label=(\roman*)]
\item one edge: $\Psi=4$;
\item a matching of $k$ edges: $\Psi=4k$;
\item a clique $K_n$: $\Psi=2n$;
\item a complete bipartite graph $K_{a,b}$:
\[
\Psi=(\sqrt a+\sqrt b)^2.
\]
In particular, a $k$-leaf star has $\Psi=(1+\sqrt{k})^2$.
\end{enumerate}
\end{proposition}

\begin{proof}
The edge case follows by symmetry.  Disjoint components add, giving the matching.
For a clique, convexity and vertex symmetry imply an optimal symmetric solution,
so $q_v=2$ for all $v$.  For $K_{a,b}$, symmetry within each side gives variables
$x,y$ with $1/x+1/y=1$ and objective $ax+by$.  Minimization yields
$x=1+\sqrt{b/a}$ and $y=1+\sqrt{a/b}$, giving the stated value.
\end{proof}

\begin{corollary}[Homogeneous scaling]
If every vertex variance is multiplied by $c>0$, then
\[
\Psi(H,c\nu^2)=c\Psi(H,\nu^2).
\]
\end{corollary}
